\documentclass[11pt]{article}
\usepackage[T1]{fontenc}
\usepackage[margin=1in]{geometry}
\usepackage{amsmath,amssymb,amsthm,booktabs}
\usepackage[hidelinks,breaklinks]{hyperref}
\hypersetup{
pdftitle={Reflection-symmetric diagonal synthesis: resolving a conjecture of Huang, Kosugi, Nishi and Matsushita},
pdfauthor={Haobo Ma; Wenlin Zhang}}
\newtheorem{theorem}{Theorem}
\newtheorem{lemma}{Lemma}
\newtheorem{corollary}{Corollary}
\newcommand{\F}{\mathbb F_2}
\newcommand{\CNOT}{\operatorname{CNOT}}
\newcommand{\diag}{\operatorname{diag}}
\newcommand{\wt}{\operatorname{wt}}
\title{Reflection-symmetric diagonal synthesis:\newline
resolving a conjecture of Huang, Kosugi,\newline Nishi and Matsushita}
\author{Haobo Ma\thanks{ChronoAI Pte Ltd, Singapore.}\and
Wenlin Zhang\thanks{National University of Singapore; ChronoAI Pte Ltd, Singapore.}}
\date{1 November 2026\\[0.4em]
\small Licensed under \href{https://creativecommons.org/licenses/by/4.0/}{CC BY 4.0}}
\begin{document}
\maketitle
\begin{abstract}
Combined with the upper bound of Huang, Kosugi, Nishi and Matsushita
[J.\ Phys.\ Soc.\ Jpn.\ \textbf{93}, 054002 (2024);
arXiv:2310.06676v2], this resolves the conjecture stated in their
Sect.~6 that $2^{2n-1}+2n-2$ CNOTs are optimal in the worst case.
For every $n\geq1$, an explicit reflection-symmetric diagonal unitary on
$2n$ qubits requires at least $2^{2n-1}+2n-2$ CNOT gates in any exact,
ancilla-free circuit over $\{\CNOT,R_z\}$, even up to global phase.
The hard instance has diagonal entry $\mathrm i$ at the all-zero and
all-one strings and entry $1$ elsewhere. Multiplicative character
invariants force all nonzero even parities to occur on rotated wires.
A birth-and-return argument then gives the additional linear term in
the CNOT bound. We give a self-contained exposition of the proof in the
Trureturing source at the cited revision.
\end{abstract}

\section{The synthesis question}
How many CNOT gates are needed in the worst case to synthesize a
reflection-symmetric diagonal unitary? Huang et al.~\cite{HKNM}
construct circuits on $2n$ wires using $2^{2n-1}+2n-2$ CNOTs and
conjecture (Sect.~6) that this count is worst-case optimal.
The explicit hard instance below
establishes the matching lower bound, by forcing even wire parities
and counting the gates needed to restore the wires.

Put $m=2n$ and index the computational basis by $x\in\F^m$.
Write $\bar x=x+\mathbf1$. A real phase function is reflection-symmetric
when $\theta(\bar x)=\theta(x)$ for every $x$.
This is exactly the two-register symmetry in Eq.~(1) of
Huang et al.~\cite{HKNM}: complementing each $n$-bit register index
complements all $m$ bits. Throughout,
\[
 R_z(\varphi)=\diag(e^{-\mathrm i\varphi/2},e^{\mathrm i\varphi/2}),
 \qquad \varphi\in\mathbb R.
\]
The allowed circuits are finite, act on these $m$ wires only, and use
distinct target and control wires for each CNOT. No ancillary qubits,
measurements, other gates, or approximation are allowed.

\begin{theorem}\label{thm:main}
For each $n\geq1$, define $\theta_*:\F^{2n}\to\mathbb R$ by
\[
 \theta_*(x)=
 \begin{cases}\pi/2,&x=0\text{ or }x=\mathbf1,\\0,&\text{otherwise},\end{cases}
 \qquad D_*=\diag(e^{\mathrm i\theta_*(x)}).
\]
Then $\theta_*(\bar x)=\theta_*(x)$. Every allowed circuit with matrix
$U=cD_*$, for any complex scalar $c$, contains at least
\[
 N_m=2^{m-1}+m-2=2^{2n-1}+2n-2
\]
CNOT gates.
\end{theorem}

The scalar is automatically nonzero, indeed of unit modulus, because
$U$ and $D_*$ are unitary. Theorem~\ref{thm:main} is existential with
respect to the class of symmetric targets: it does not say that every
member of that class is hard. The identity is an immediate easy member.
The lower-bound statement and proof are already present in the
Trureturing formalization~\cite{upstream}.

Parity networks give a standard description of CNOT/phase circuits.
Amy, Azimzadeh and Mosca~\cite{AAM}, Sect.~2, define these networks
with a specified final linear map. Identity parity networks return
that map to the identity.

Cao et al.~\cite{CLNST} study the graphic case, where the required
labels have weight two. For a connected graph their lower bound is
$|E|+|V|-1$. It gives $9$ for the complete graph on four wires.
Here the weight-four label is also forced, and the lower bound is $10$.
The argument below treats every nonzero even label and controls
the freedom to change real phases by multiples of $2\pi$.

\section{Wire labels and multiplicative characters}
Represent a wire's current Boolean value as $w\cdot x$, with label
$w\in\F^m$. Initially the labels are the standard basis vectors $e_j$.
A CNOT with target $t$ and control $c$ replaces $w_t$ by $w_t+w_c$.
Thus the label matrix stays invertible and its rows never vanish.
A rotation on a wire labelled $w$ multiplies the basis amplitude by
\[
 g(w\cdot x),\qquad g(0)=e^{-\mathrm i\varphi/2},\quad
 g(1)=e^{\mathrm i\varphi/2}.
\]
Induction over the circuit consequently gives a monomial normal form
\[
 |x\rangle\longmapsto f(x)|Ax\rangle,\qquad
 f(x)=\prod_k g_k(w_k\cdot x),\qquad A\in\mathrm{GL}(m,\F).
\]
All factors are nonzero. If the circuit is diagonal up to a scalar,
its nonzero entry on each input cannot change the basis index, so
$Ax=x$ for every $x$. Hence $A=I$ and the final wire labels are again
$e_1,\ldots,e_m$.

For any nonzero-valued function $h:\F^m\to\mathbb C$ and any label $v$,
define
\[
 \mathcal I_v(h)=\prod_{x\in\F^m}h(x)^{(-1)^{v\cdot x}}.
\]
The exponents are integers $+1$ or $-1$, so negative exponents mean
multiplicative inverses, not a choice of logarithm or phase branch.
In particular, $\mathcal I_v(h_1h_2)=\mathcal I_v(h_1)\mathcal I_v(h_2)$.

\begin{lemma}\label{lem:character}
If $v\ne0$, then $\mathcal I_v(c)=1$ for each constant $c\ne0$.
If also $w\ne0$ and $w\ne v$, then
$\mathcal I_v(x\mapsto g(w\cdot x))=1$ for every nonzero-valued
$g:\F\to\mathbb C$.
\end{lemma}
\begin{proof}
Choose $u$ with $v\cdot u=1$. Pairing $x$ with $x+u$ cancels the
constant factors in the first assertion. For the second assertion,
the distinct nonzero vectors $v,w$ are linearly independent over $\F$.
Therefore the map $x\mapsto(v\cdot x,w\cdot x)$ is onto $\F^2$,
and there is a $u$ with $v\cdot u=1$ and $w\cdot u=0$.
The same fixed-point-free pairing preserves $g(w\cdot x)$ and reverses
its exponent. Each pair contributes $1$.
\end{proof}

Call $v$ even when $\wt(v)$ is even, equivalently $v\cdot\mathbf1=0$.
The diagonal function $d_*(x)=e^{\mathrm i\theta_*(x)}$ is $\mathrm i$
at $0,\mathbf1$ and $1$ elsewhere. For every nonzero even $v$,
\[
 \mathcal I_v(d_*)=\mathrm i\,\mathrm i=-1,
\]
since both marked strings have dot product zero with $v$.
If $v$ never occurred on a rotated wire, Lemma~\ref{lem:character}
would give $\mathcal I_v(f)=1$. But $f=cd_*$ for a realizing circuit,
so it also gives $\mathcal I_v(f)=\mathcal I_v(c)\mathcal I_v(d_*)=-1$,
a contradiction. Thus \emph{every nonzero even label must occur on a
rotated wire}. This argument is insensitive to phase wrapping.

\section{Counting even outputs and returns}
Let $H=\{v\in\F^m:v\cdot\mathbf1=0\}$. It has dimension $m-1$,
and contains $2^{m-1}-1$ nonzero vectors. All initial wire labels are
odd. The first appearance of any even label must therefore be a CNOT
output; rotations do not change labels. Since a single CNOT produces
one output label, at least $2^{m-1}-1$ CNOTs have even output.

An update whose target and control labels are both odd is a
\emph{birth}; its output is even. At every stage, every even label ever
produced lies in the span of birth outputs produced so far. To prove
this invariant, initially there are no even labels. An odd--odd update
adds its birth output as a generator. An even--even update outputs the
sum of two labels already in the span. The two mixed-parity updates
produce odd labels, so they add no even label. Unchanged rows preserve
the invariant. This proves the assertion inductively, including labels
that later disappear from the current rows.

All nonzero vectors of $H$ must be produced, so the birth outputs span
$H$. If $B$ is the number of births, this forces $B\geq m-1$.
Now count the number $E$ of even wire labels currently present:
\begin{center}
\begin{tabular}{llll}
\toprule
Target before & Control & Output & Change in $E$\\
\midrule
odd & odd & even & $+1$ (birth)\\
even & odd & odd & $-1$ (return)\\
even & even & even & $0$\\
odd & even & odd & $0$\\
\bottomrule
\end{tabular}
\end{center}
The initial and final labels are standard basis rows, so $E$ starts
and ends at zero. Hence the number $R$ of returns equals $B$, and
$R\geq m-1$. Returns have odd output and are disjoint from the
even-output CNOTs already counted. Therefore the total CNOT count is
\[
 K\geq (2^{m-1}-1)+R\geq(2^{m-1}-1)+(m-1)=N_m.
\]
Together with the complement invariance of $\theta_*$, this proves
Theorem~\ref{thm:main}.

\begin{corollary}[Using the literature upper bound]
In the stated model, the worst-case optimal CNOT count among all
reflection-symmetric diagonal unitaries on $2n$ wires is $N_{2n}$.
\end{corollary}
\begin{proof}
Algorithm 1 of Huang et al.~\cite{HKNM}, Sect.~4 and Appendix C,
implements every target in the class up to global phase with
$(2^{2n-1}-1)+(2n-1)=N_{2n}$ CNOTs. This is the literature upper-bound
premise. Theorem~\ref{thm:main} supplies a target meeting it from below.
\end{proof}

\paragraph{Higher-dimensional count and remaining questions.}
For integers $d,n\geq1$, substituting $dn$ for $n$ gives the exact
worst-case count $2^{2dn-1}+2dn-2$ on $2dn$ wires. In particular,
the $6n$-wire count $2^{6n-1}+6n-2$ in Sect.~6 of Huang et al.\ is
sharp for the full class of complement-symmetric diagonal targets.
Their three-dimensional interaction-potential family, defined by
Eq.~(5), has additional radial symmetry: the lower bound here does
not establish worst-case optimality within that narrower family.
Section~6 also asks for optimal circuit depth. The gate-count proof
does not determine that depth, since disjoint gates can run in parallel.
Thus the depth and restricted-potential synthesis questions remain
unresolved by this argument.

\section{A finite synthesis check}
With $\chi_v(x)=(-1)^{v\cdot x}$,
\[
 \theta_*(x)=\frac{\pi}{2^m}\sum_{v\in H}\chi_v(x).
\]
This follows because a character sum over $H$ is $|H|$ on
$H^\perp=\{0,\mathbf1\}$ and zero elsewhere: outside $H^\perp$,
translation of the summation index by a vector pairing nontrivially
with $x$ cancels terms. Rotating once at each nonzero even label by
$-\pi/2^{m-1}$ thus implements $e^{-\mathrm i\pi/2^m}D_*$ whenever
the labels return to identity.

Exact breadth-first searches over ordered wire labels and visited
nonzero even labels give shortest identity-return networks of lengths
$2$ and $10$ for $m=2$ and $m=4$, respectively. Inserting these rotations
at each label's first occurrence gives circuits for $D_*$ of those
lengths. Integer and rational-arithmetic checks also verify every
hard-instance character and every basis phase in the displayed identity
for $m=2,4,6,8$. These are finite checks, distinct from the all-$n$
written proof in Sections~2--3. The pinned declaration
\texttt{ReflectionDiagonalCNOTLowerBound.result} proves the lower bound;
worst-case equality additionally uses the cited synthesis upper bound.
The character and phase checks cover $m=2,4,6,8$, and the shortest-network
checks cover $m=2,4$. The exact checker and witnesses accompany
the source package.

\begin{thebibliography}{9}
\bibitem{HKNM}
X.~Huang, T.~Kosugi, H.~Nishi and Y.-i.~Matsushita,
\emph{Optimized Synthesis of Circuits for Diagonal Unitary Matrices
with Reflection Symmetry}, Journal of the Physical Society of Japan
\textbf{93}(5), 054002 (2024).
\href{https://doi.org/10.7566/JPSJ.93.054002}{doi:10.7566/JPSJ.93.054002}.
\href{https://arxiv.org/abs/2310.06676v2}{arXiv:2310.06676v2},
Eq.~(1), Algorithm 1, Sect.~6 and Appendix C.
\bibitem{AAM}
M.~Amy, P.~Azimzadeh and M.~Mosca,
\emph{On the CNOT-complexity of CNOT-PHASE circuits},
published as \emph{On the controlled-NOT complexity of controlled-NOT--phase
circuits}, Quantum Science and Technology \textbf{4}, 015002 (2019),
\href{https://doi.org/10.1088/2058-9565/aad8ca}{doi:10.1088/2058-9565/aad8ca};
\href{https://arxiv.org/abs/1712.01859v2}{arXiv:1712.01859v2}
(2018), Sect.~2.
\bibitem{CLNST}
Y.~Cao, Y.~Lu, J.~Nie, X.~Sun and G.~Tian,
\emph{Toward Minimum Graphic Parity Networks},
\href{https://arxiv.org/abs/2509.10070v1}{arXiv:2509.10070v1}
(2025), Theorem~1.1.
\bibitem{upstream}
Trureturing repository contributors, version-pinned sources:
\href{https://github.com/the-omega-institute/trureturing/blob/02c695f8c17e35965f80603427074b1dd434eed9/D5/S3/Quantum/Magic/ReflectionDiagonalCNOTLowerBound.lean}{ReflectionDiagonalCNOTLowerBound.lean}
and
\href{https://github.com/the-omega-institute/trureturing/blob/02c695f8c17e35965f80603427074b1dd434eed9/D5/S3/Quantum/Magic/ReflectionDiagonalParityCounting.lean}{ReflectionDiagonalParityCounting.lean}.
The lower bound is \texttt{result : claim} in the first module;
the second supplies \texttt{count\_bound}.
\end{thebibliography}
\end{document}
